\section{Model}
\label{sec:model}

\subsection{Parties and Orders}

The system has $n$ \emph{officers} and a set of \emph{devices}. An officer is a person with
a console. A device does orders: for example, it opens a breaker. One officer at a time is
the \emph{commander}. A second officer can be the \emph{deputy} of the commander.

Each officer has a key and can sign messages. Each party knows the public key of each
officer. An \emph{order} is a message that an officer signs. A device \emph{accepts} an
order when it does the order.

\subsection{Attacker}

The system has one attacker. The model gives the attacker these properties.
\begin{description}
  \item[A1] \emph{Links.} The attacker controls all links. It reads each message. It can delay a
  message without a time limit, stop it, replay it, and change the sequence of messages.
  \item[A2] \emph{Keys.} The attacker can get the key and the data of an officer at a time in a run.
  From that time, the attacker can sign with that key. An officer is \emph{hostile} in a
  run if the attacker knows the key of the officer at some time in the run. The other
  officers are \emph{loyal}. Not more than $f$ officers are hostile in a run, and
  $f \ge 1$.
  \item[A3] \emph{Limits.} The attacker cannot forge the signature of a loyal officer. The attacker
  cannot change the clock or the data of a device. The attacker is not at the location of
  a device.
\end{description}

A hostile officer can be a person who helps the attacker. A hostile officer can also be a
loyal person with a console that the attacker controls. The model does not make a
difference between these two conditions.

\subsection{Time, Clocks, Runs and Views}

Real time is a number $t \ge 0$. No party knows real time. Each device and each loyal
officer $p$ has a clock $C_p$, and $|C_p(t) - t| \le \eps$ for all $t$. All parties know
the clock error $\eps$. \Cref{sec:model-limits} examines attacks on clocks.

A \emph{run} is one possible sequence of the steps of all parties, with the real time of
each step. Loyal officers and devices follow the protocol. A device is deterministic, and
messages are its only input. \Cref{sec:protocol} adds one more input, a switch at the
device.

The \emph{view} of a device at a time $t$ has two parts. The first part is the data that
the device has at the start. The second part is each message that the device received
until $t$, with the clock time at which the device received it. The steps of a device are
a function of its view. Thus, if a device has the same view in two runs, it accepts the
same orders at the same times.

\subsection{Quantities}
\label{sec:quantities}

The definitions that follow apply to each protocol in this model. They use only the orders
that a device accepts. Thus they do not use the rules of one protocol.

\begin{dfn}[Outage tolerance]
\label{def:tolerance}
Let $d$ be a device and $P \ge 0$. A protocol has \emph{outage tolerance} $P$ at $d$ if it
has a run with these properties. All officers are loyal. From a time $t_0$, the commander
$c$ and the device $d$ receive only messages that $c$ or $d$ sent. The device $d$ accepts an
order of $c$ at a time $t_1 \ge t_0 + P$.
\end{dfn}

\begin{dfn}[Hostile tenure]
\label{def:tenure}
Let $d$ be a device and $\Theta \ge 0$. A protocol has \emph{hostile tenure} $\Theta$ at
$d$ if it has a run with these properties. At a time $t_0$, the attacker gets the key of
the commander $c$, and it gets no other key in the run. From $t_0$, each other officer does
the steps of the protocol that remove $c$. The device $d$ accepts an order of $c$ at a time
$t_1 \ge t_0 + \Theta$.
\end{dfn}

A large outage tolerance is good, and a large hostile tenure is bad. The two definitions
give times, but a protocol can also put a limit on the number of orders. \Cref{sec:limits}
includes that condition.

\begin{dfn}[Deputy period]
\label{def:deputy}
Let $d$ be a device and $F \ge 0$. A protocol has \emph{deputy period} $F$ at $d$ if it has
a run with these properties. All officers are loyal. From a time $t_0$, the deputy $b$ and
the device $d$ receive only messages that $b$ or $d$ sent. The device $d$ accepts an order
of $b$ at a time $t_1 \ge t_0$ and at a time $t_2 \ge t_1 + F$.
\end{dfn}

In the run for the deputy period, the device gets no message from the commander. The device
cannot know the cause. It is possible that the commander stopped. It is also possible that
the links to the commander have a failure, or that an attacker cut them.

\begin{dfn}[Deputy exposure]
\label{def:exposure}
Let $d$ be a device and $F \ge 0$. A protocol has \emph{deputy exposure} $F$ at $d$ if it
has a run with these properties. At a time $t_0$, the attacker gets the key of the deputy
$b$, and it gets no other key in the run. The commander is loyal, and the attacker stops no
message between the commander and the other loyal officers. The device $d$ accepts an order
of $b$ at a time $t_1 \ge t_0$ and at a time $t_2 \ge t_1 + F$.
\end{dfn}

A protocol that has one of these four properties for a value also has it for each smaller
value. When we give one number as the outage tolerance of a protocol, it is the smallest
number that is not less than each such value. The same applies to the other three
quantities.

\begin{dfn}[Unity of command]
\label{def:unity}
Let $\gamma > 0$. A protocol has \emph{unity of command with separation} $\gamma$ if each
run has this property. Let a device accept an order of an officer $x$ at real time $t$.
Let a device accept an order of a different officer $y$ at real time $t'$. Then
$|t' - t| \ge \gamma$.
\end{dfn}
